\subsection{Checks requested in review}\label{sec:res-review}
% src: S-PAPER-12, S-PAPER-8
Reviewers asked for further checks.
We ran them in one run, \RThreeRun.
We fixed the rules before the run in \texttt{analysis/prereg\_review\_fixes.md}.
All checks are descriptive and add no verdict.
Each new row sits next to the original Q1 row and does not replace it.
The pre-registered Q1 verdict and its sentence stay as they are.
Every check uses the \QOneN{} kept markets, except where stated.

% src: S-PAPER-12
Check M1 asks whether the price definition drives the net gap.
Under the volume-weighted mean price (VWAP), the net gap is \RThreeMOneVwapNetGap{} (\CiLevel{} CI [\RThreeMOneVwapNetGapCiLow, \RThreeMOneVwapNetGapCiHigh]).
Under the trade-based midpoint, it is \RThreeMOneMidNetGap{} (\CiLevel{} CI [\RThreeMOneMidNetGapCiLow, \RThreeMOneMidNetGapCiHigh]).
The original row has \GapNetDiffMain.
Both intervals lie below 0.
So the pre-registered reading is that the price rule does not create the gap.
The midpoint is a trade-based proxy and not a quote midpoint.
It exists for the \RThreeMOneMidN{} markets with both trade sides (\RThreeMOneShareBoth).
The mean half spread of this proxy is $h=\RThreeMOneMeanHalfSpread$.
Is the size of the main net gap within $h$? \RThreeMOneGapWithinHalfSpread.
\Cref{fig:r3-m1} shows the three net gaps.

\begin{figure}[H]
\centering
\figorbox{figures/fig_r3_m1_price_definition}{0.8\linewidth}
\caption{Net gap under each price definition, with \CiLevel{} intervals. Why: an interval below 0 under VWAP and the midpoint means the gap is not a price-rule artifact. Source: runs \RThreeRun{} and \RThreeFigRun.}
\label{fig:r3-m1}
\end{figure}

% src: S-PAPER-12, S-PAPER-8
Check M2 splits the markets by label.
In the markets labelled Yes/No ($n=\RThreeMTwoYesNoN$), Yes is priced above its observed rate.
The net gap is \RThreeMTwoYesNoNetGap{} (\CiLevel{} CI [\RThreeMTwoYesNoNetGapCiLow, \RThreeMTwoYesNoNetGapCiHigh]).
Of \NBins{} bins, \RThreeMTwoYesNoBinsBh{} are off after BH.
The other markets show no rejection ($n=\RThreeMTwoOtherN$, $\chi^2=\RThreeMTwoOtherChiSq$, $p=\RThreeMTwoOtherP$).
Elsewhere we keep the term outcome 1.
\Cref{fig:r3-m2} shows the markets labelled Yes/No.
The figure uses Wilson intervals, which are outside the lectures.

\begin{figure}[H]
\centering
\figorbox{figures/fig_r3_m2_yes_reliability}{0.55\linewidth}
\caption{Observed Yes rate against mean Yes price in the markets labelled Yes/No. Why: points below the diagonal mean Yes is priced above its observed rate. Source: runs \RThreeRun{} and \RThreeFigRun.}
\label{fig:r3-m2}
\end{figure}

% src: S-PAPER-12
Check M3 asks how much dependence between markets matters.
The data hold no event key, so we use two proxy keys.
The reference row k0 treats each market alone, with no cluster, and is not a key.
The first key is \texttt{\RThreeMThreeReadProxy}.
It takes the first three tokens of the market slug.
It has \RThreeMThreeKOneNClusters{} clusters.
The second key is category by close hour, with \RThreeMThreeKTwoNClusters{} clusters.
We read the first key because its net gap interval is wider.
Under it, the \CiLevel{} interval is [\RThreeMThreeKOneNetGapCiLow, \RThreeMThreeKOneNetGapCiHigh], and it still lies below 0.
Under the second key it is [\RThreeMThreeKTwoNetGapCiLow, \RThreeMThreeKTwoNetGapCiHigh].
For the first key the cluster chi-square is \RThreeMThreeKOneChiSq{} ($p=\RThreeMThreeKOneP$).
For the second key it is \RThreeMThreeKTwoChiSq{} ($p=\RThreeMThreeKTwoP$).
The cluster bootstrap chi-square is outside the lectures.
The proxies are not true event keys.
\Cref{fig:r3-m3} compares the intervals.

\begin{figure}[H]
\centering
\figorbox{figures/fig_r3_m3_cluster_bootstrap}{0.8\linewidth}
\caption{Net gap with a cluster bootstrap interval under each proxy key, next to the one-market-per-cluster row. Why: a wider interval means markets in a cluster move together. Source: runs \RThreeRun{} and \RThreeFigRun.}
\label{fig:r3-m3}
\end{figure}

% src: S-PAPER-12
Check M4 compares the kept markets with the dropped ones.
The kept markets have a median volume before $T-\DaysShort$ day of \RThreeMFourKeptMedianVolumeBeforeDay{} USDC.
The pooled eligible dropped markets have \RThreeMFourDroppedMedianVolumeBeforeDay{} USDC.
The median duration is \RThreeMFourKeptMedianDuration{} against \RThreeMFourDroppedMedianDuration{} days.
The median trade count is \RThreeMFourKeptMedianTrades{} against \RThreeMFourDroppedMedianTrades.
The stale-price markets have a net gap of \RThreeMFourStaleNetGap{} (\CiLevel{} CI [\RThreeMFourStaleNetGapCiLow, \RThreeMFourStaleNetGapCiHigh]).
\Cref{fig:r3-m4} shows the share of markets in each group.

\begin{figure}[H]
\centering
\figorbox{figures/fig_r3_m4_kept_vs_dropped}{0.8\linewidth}
\caption{Share of in-scope markets in each kept and dropped group. Why: it shows how much of the market set each drop rule removes. Source: runs \RThreeRun{} and \RThreeFigRun.}
\label{fig:r3-m4}
\end{figure}

% src: S-PAPER-12
Check M5 asks how large the miscalibration is in usable terms.
We fit a logistic calibration of the outcome on the price, $\mathrm{logit}\,P(y=1)=a+b\,\mathrm{logit}(p)$.
The intercept is $a=\RThreeMFiveA$ (Wald \CiLevel{} CI [\RThreeMFiveAWaldLow, \RThreeMFiveAWaldHigh]).
The slope is $b=\RThreeMFiveB$ (Wald \CiLevel{} CI [\RThreeMFiveBWaldLow, \RThreeMFiveBWaldHigh]).
The practical line is the mean half spread, \RThreeMOneMeanHalfSpread.
The largest implied gap on the price grid is \RThreeMFiveMaxAbsDelta{} in size.
Of \RThreeMFiveGridPrices{} grid prices, \RThreeMFivePracticalPrices{} pass the practical line.
The Brier skill score against climatology is \RThreeMFiveSkill{} (Brier \RThreeMFiveBrier{} against \RThreeMFiveClimatologyBrier).
The climatology is in sample.
The majority-class accuracy is \RThreeMFiveMajorityAccuracy.
The accuracy at the \FavoriteLine{} line is \RThreeMFiveAccuracyHalf.
Logistic calibration and the skill score are outside the lectures.
\Cref{fig:r3-m5} shows the implied gap.

\begin{figure}[H]
\centering
\figorbox{figures/fig_r3_m5_implied_gap}{0.8\linewidth}
\caption{Implied gap between the fitted probability and the price, against the practical line. Why: a band that excludes 0 at a price shows the gap there is not noise. Source: runs \RThreeRun{} and \RThreeFigRun.}
\label{fig:r3-m5}
\end{figure}

% src: S-PAPER-12
The duplicate-row check removes exact duplicate trade rows.
They are \RThreeMinorFourDupShare{} of the trade rows of the kept markets.
Removing them drops \RThreeMinorFourRemoved{} markets from the kept set.
Then $\chi^2=\RThreeMinorFourChiSq$ against \QOneChiSq{} in the original row.
The net gap is \RThreeMinorFourNetGap{} against \GapNetDiffMain.
The top \TopOneLabel{} maker share changes by \RThreeMinorFourMakerTopOneDeltaPp{} percentage points.
The top \TopTwentyLabel{} maker share changes by \RThreeMinorFourMakerTopTwentyDeltaPp.
For takers the changes are \RThreeMinorFourTakerTopOneDeltaPp{} and \RThreeMinorFourTakerTopTwentyDeltaPp.

% src: S-PAPER-12
Two peer questions ask whether the gap depends on volume and on the close period.
We split the kept markets into volume terciles, by volume before $T-\DaysShort$ day.
The cut points are \RThreePeerACutLowMid{} and \RThreePeerACutMidHigh{} USDC.
The net gap is \RThreePeerALowNetGap{} in the low tercile (CI [\RThreePeerALowNetGapCiLow, \RThreePeerALowNetGapCiHigh], BH $p=\RThreePeerALowBhP$).
It is \RThreePeerAMidNetGap{} in the middle tercile (CI [\RThreePeerAMidNetGapCiLow, \RThreePeerAMidNetGapCiHigh], BH $p=\RThreePeerAMidBhP$).
It is \RThreePeerAHighNetGap{} in the high tercile (CI [\RThreePeerAHighNetGapCiLow, \RThreePeerAHighNetGapCiHigh], BH $p=\RThreePeerAHighBhP$).
The high minus low difference is \RThreePeerADiff{} (CI [\RThreePeerADiffCiLow, \RThreePeerADiffCiHigh], $p=\RThreePeerADiffP$).
\Cref{fig:r3-peer-a} shows the terciles.

\begin{figure}[H]
\centering
\figorbox{figures/fig_r3_peer_a_volume}{0.8\linewidth}
\caption{Net gap per volume tercile, with \CiLevel{} intervals. Why: the pre-registered difference test compares the high and low terciles, with high minus low \RThreePeerADiff{} (CI [\RThreePeerADiffCiLow, \RThreePeerADiffCiHigh], $p=\RThreePeerADiffP$). Source: runs \RThreeRun{} and \RThreeFigRun.}
\label{fig:r3-peer-a}
\end{figure}

The close periods are \RThreePeerBPeriodOneLabel, \RThreePeerBPeriodTwoLabel, and \RThreePeerBPeriodThreeLabel.
The net gap is \RThreePeerBPeriodOneNetGap{} in \RThreePeerBPeriodOneLabel{} (CI [\RThreePeerBPeriodOneNetGapCiLow, \RThreePeerBPeriodOneNetGapCiHigh], BH $p=\RThreePeerBPeriodOneBhP$, $n=\RThreePeerBPeriodOneN$).
It is \RThreePeerBPeriodTwoNetGap{} in \RThreePeerBPeriodTwoLabel{} (CI [\RThreePeerBPeriodTwoNetGapCiLow, \RThreePeerBPeriodTwoNetGapCiHigh], BH $p=\RThreePeerBPeriodTwoBhP$, $n=\RThreePeerBPeriodTwoN$).
In \RThreePeerBPeriodThreeLabel{} the chi-square test does not reject (BH $p=\RThreePeerBPeriodThreeBhP$, $n=\RThreePeerBPeriodThreeN$).
In the same period the net gap is \RThreePeerBPeriodThreeNetGap.
Its \CiLevel{} CI [\RThreePeerBPeriodThreeNetGapCiLow, \RThreePeerBPeriodThreeNetGapCiHigh] excludes 0.
\Cref{fig:r3-peer-b} shows the periods.

\begin{figure}[H]
\centering
\figorbox{figures/fig_r3_peer_b_period}{0.8\linewidth}
\caption{Net gap per close period, with \CiLevel{} intervals. Why: overlapping intervals mean no sign that the gap changes with the period. Source: runs \RThreeRun{} and \RThreeFigRun.}
\label{fig:r3-peer-b}
\end{figure}
